\section{感受野决定了网络能看多远}
\label{sec:14-Deep-Learning/04-Convolutional-Networks/02}

你在做卫星图的道路提取。输入 \(512\times512\)，网络是十二层 \(3\times3\) 卷积、stride 全为 1，逐像素输出这里是不是路。直路段的准确率很高，一到被树冠遮住的路口就断成两截——被遮的那几十个像素在图上和裸土几乎一样，唯一能救它的证据是两侧延伸过来的路面走向。你把网络加宽一倍、多训了十轮，断口照旧。

问题不在容量。问题是那个输出像素根本没有机会看到两侧的路。它能看到多远是一个可以事先算出来的数，算完再决定改什么。

\subsection{递推式里真正在动的是采样间隔}

设第 \(l\) 层的核宽 \(k_l\)、步幅 \(s_l\)、空洞率 \(d_l\)。记 \(r_l\) 为第 \(l\) 层一个单元在原始输入上覆盖的宽度，\(j_l\) 为第 \(l\) 层相邻两个单元的中心在原始输入上相隔多少个像素——把它叫作跳距。两条递推：

\[
\begin{aligned}
j_l &= j_{l-1}\,s_l,\\
r_l &= r_{l-1} + d_l\,(k_l-1)\,j_{l-1},
\end{aligned}
\]

初值 \(j_0=r_0=1\)。展开跳距得 \(j_{l-1}=\prod_{i<l}s_i\)，于是

\[
r_L = 1 + \sum_{l=1}^{L} d_l\,(k_l-1)\prod_{i<l}s_i .
\]

第二条式子值得停一下。第 \(l\) 层的核横跨 \(k_l\) 个输入单元，这些单元本身在原图上彼此相隔 \(j_{l-1}\) 个像素，所以核的两端在原图上拉开了 \((k_l-1)j_{l-1}\) 的距离；空洞率 \(d_l\) 再把这个跨度乘 \(d_l\) 倍。感受野每层增加的量，等于本层核跨过的格数乘以每格在原图上值多少像素。stride 出现在乘积里而不是加法里，因为它改变的是后续每一层每一格的兑换率——一个 stride 2 会让它之后所有层的每一步都值两个原始像素。

回到那个道路网络。全是 \(k=3,s=1,d=1\)，每层加 2，十二层的 \(r_{12}=25\)。\(512\times512\) 的图上，一个 25 像素宽的窗口，遮挡稍宽一点就整个落在窗口外面。加宽通道解决不了这件事：宽度改变的是能提取多少种特征，不改变能取到哪些像素。

\subsection{小核堆叠为什么划算}

要把 \(r\) 拉到 25，也可以只用一层 \(25\times25\) 的卷积。同样的感受野，两种做法的账完全不同。

参数上，\(C\) 个输入通道、\(C\) 个输出通道时，单层大核是 \(C^2\cdot 25^2=625\,C^2\)，十二层小核是 \(12\,C^2\cdot 3^2=108\,C^2\)，差近六倍。感受野对深度是线性增长而对单层核宽也是线性增长，可参数量对核宽是平方增长——同样一份视野，摊到多层去买要便宜得多。二维下这个比值随目标感受野继续拉大。

比参数更要紧的是非线性次数。单层大核在整个 \(25\times25\) 窗口上只做一次线性组合再过一次激活；十二层小核在同一片区域里插了十二次激活。中间那些层可以先判断``这里有没有路面纹理''，再在下一层基于这个判断去看邻域，而不是被迫用一次线性投影同时完成检测与聚合。堆叠买到的不只是视野，还有视野内部的分层决策。

这笔账并非在所有情形都倒向小核。要达到 \(r\) 的深度是 \(O(r)\) 层，每层都带来一次梯度传递和一次激活的显存占用；当目标感受野很大时，纯靠 \(3\times3\) 堆出来的深度本身会变成训练与推理的负担。近年重新起用 \(7\times7\) 乃至更大核的架构，通常配合 depthwise 让参数量不再是平方级的那一项，用较浅的网络换掉一部分深度。小核堆叠的优势成立于``参数按 \(C^2k^2\) 计价''这个前提，前提变了结论也要重算。

\subsection{空洞卷积把线性增长换成指数增长}

递推式里 \(d_l\) 是乘在增量上的，参数量却和 \(d_l\) 无关——核仍然只有 \(k^2\) 个数，只是采样点之间被拉开了。逐层令 \(d_l=2^{l-1}\)、\(k=3\)、\(s=1\)，则每层增量为 \(2\cdot 2^{l-1}\)，累加得

\[
r_L = 1 + \sum_{l=1}^{L} 2^{l} = 2^{L+1}-1 .
\]

四层就到 31，八层到 511，覆盖整张图。同样的 \(r=25\)，普通堆叠要十二层，空洞堆叠要四层。

代价写在采样模式里。\(d=2\) 的核只取偶数偏移的位置，紧邻的那些像素被跳过；如果连续几层的空洞率有公因子，被跳过的位置会在多层之间对齐，输出像素之间形成互不相交的采样格子——同一个输出只用到了输入的一个稀疏子格，相邻输出用的是另一个子格，彼此没有交换过信息。这就是常说的网格伪影，表现为分割边界上的条纹。规避办法是让相邻层的空洞率互素，比如 \(1,2,3\) 循环而不是 \(1,2,4\)，让不同层的采样格互相填补。感受野的宽度和感受野内部的密度是两件事，递推式只管前者。

\subsection{理论感受野不等于有效感受野}

\begin{center}
\begin{tikzpicture}[x=0.62cm,y=0.85cm,>=stealth]
% 输入行与三层特征行
\foreach \i in {0,...,14} {
  \draw[gray!45] (\i,0) rectangle +(0.86,0.5);
}
\foreach \i in {4,...,10} {
  \fill[gray!30] (\i,0) rectangle +(0.86,0.5);
  \draw[gray!60] (\i,0) rectangle +(0.86,0.5);
}
\foreach \i in {5,...,9} {
  \fill[gray!30] (\i,1) rectangle +(0.86,0.5);
  \draw[gray!60] (\i,1) rectangle +(0.86,0.5);
}
\foreach \i in {6,...,8} {
  \fill[gray!30] (\i,2) rectangle +(0.86,0.5);
  \draw[gray!60] (\i,2) rectangle +(0.86,0.5);
}
\fill[gray!55] (7,3) rectangle +(0.86,0.5);
\draw[gray!80] (7,3) rectangle +(0.86,0.5);
% 锥形边界
\draw[dashed] (7.43,3) -- (4.0,0.5);
\draw[dashed] (7.43,3) -- (11.3,0.5);
\node[right,font=\scriptsize] at (15.2,3.25) {\(r_3=7\)};
\node[right,font=\scriptsize] at (15.2,2.25) {\(r_2=5\)};
\node[right,font=\scriptsize] at (15.2,1.25) {\(r_1=3\)};
\node[right,font=\scriptsize] at (15.2,0.25) {输入};
% 有效感受野权重曲线
\draw[->] (3.0,-2.1) -- (12.2,-2.1) node[right,font=\scriptsize] {输入位置};
\draw[very thick] plot[domain=3.6:11.3,samples=70]
  (\x,{-2.1+1.35*exp(-(\x-7.43)*(\x-7.43)/2.2)});
\draw[dotted] (4.0,-2.1) -- (4.0,0.5);
\draw[dotted] (11.3,-2.1) -- (11.3,0.5);
\draw[dotted] (6.38,-2.1) -- (6.38,-1.28);
\draw[dotted] (8.48,-2.1) -- (8.48,-1.28);
\draw[<->] (6.38,-2.45) -- (8.48,-2.45);
\node[below,font=\scriptsize] at (7.43,-2.45) {有效感受野};
\draw[<->] (4.0,-3.15) -- (11.3,-3.15);
\node[below,font=\scriptsize] at (7.43,-3.15) {理论感受野};
\end{tikzpicture}
\end{center}

\emph{三层 \(3\times3\) 堆出的锥形覆盖七个输入位置，但这七个位置对输出的贡献按近似 Gaussian 衰减，中心几格几乎决定了全部。}

图里锥形是理论感受野：从输出单元逆着连接往回追，能追到的输入位置就是它。这个集合只回答``有没有一条路径''，不回答``这条路径传了多少''。

看权重就知道差别有多大。仍取 \(k=3,s=1\) 的 \(L\) 层堆叠，暂时去掉激活，那么输出对输入位置的偏导是所有路径贡献之和。每经过一层，路径可以把偏移量改变 \(-1,0,+1\) 三者之一，走完 \(L\) 层的总偏移是这 \(L\) 个独立选择之和。抵达偏移 \(\delta\) 的路径条数，就是均匀分布在 \(\{-1,0,+1\}\) 上的随机变量做 \(L\) 次卷积后落在 \(\delta\) 处的质量，中心极限定理说它趋于 Gaussian，方差为 \(L\cdot\Var(\text{单层偏移})=\tfrac{2}{3}L\)。于是贡献的有效宽度按 \(\sqrt{L}\) 增长，而理论感受野半径按 \(L\) 增长，比值

\[
\frac{\text{有效半径}}{\text{理论半径}}\;\sim\;\sqrt{\frac{2}{3L}}\;\xrightarrow[L\to\infty]{}\;0 .
\]

条件对账要摆清楚。这条推理只数了路径条数，等于假设各条路径的权重同号且量级相当；真实权重有正有负，路径之间会相消，实测的有效感受野只会比这个估计更小。ReLU 让一部分路径在给定输入下彻底断掉，进一步削减有效宽度，同时使它随输入而变——有效感受野不是架构的固定属性，而是依赖当前样本的量。残差连接给偏移量分布加了一个集中在 0 的分量，反而让衰减更陡；空洞卷积则把单层偏移的方差从 \(\tfrac{2}{3}\) 抬到 \(\tfrac{2}{3}d^2\)，这也是它比堆深度更有效的另一半原因。

结论落回工程。\(r_L\) 覆盖全图，从来不代表网络在用全图。一个名义感受野 \(500\) 像素的分类网络，实际决策仍可能主要由目标周围几十个像素的纹理决定；它在纹理与形状冲突的样本上偏向纹理，在需要跨越大片遮挡的样本上失灵。要让远处的证据真的进来，得改变贡献的分布本身：用空洞把偏移方差抬上去，用下采样让每一步在原图上值更多像素，或者干脆换成不受路径衰减约束的全局读取机制——那是\hyperref[sec:14-Deep-Learning/05-Attention-and-Sequence-Models/01]{``Self-Attention 是内容决定的加权读取''}要做的事。

回到开头的道路网络。十二层 \(3\times3\) 的 \(r_{12}=25\) 已经不够，而有效宽度大约只有 \(\sqrt{2\cdot12/3}\approx 2.8\) 个像素的量级——它其实一直在靠路面的局部纹理判断，遮挡一来就没有可依据的东西。把中段几层换成空洞率 \(1,2,3,1,2,3\) 的循环，理论感受野和有效宽度会一起上去，这比继续加宽通道更对症。

下一篇处理另一条把感受野推远的路径：下采样。它让每一步在原图上值更多像素，代价却不是参数或深度，而是空间分辨率本身。
